\chapter{Planteamiento del Problema}

En la Universidad Nacional, los procesos de gestión de evidencias para la acreditación de carreras presentan dificultades asociadas con la dispersión, duplicidad y falta de trazabilidad de los documentos institucionales. Estas limitaciones provocan retrasos en la preparación de informes para los procesos de acreditación del SINAES y aumentan el riesgo de incumplimiento de los estándares de calidad exigidos.

De forma paralela, la administración de los tiempos de jornada universitarios se lleva a cabo mediante procedimientos manuales o apoyados en hojas de cálculo, lo cual genera inconsistencias en el cálculo de cargas docentes, dificultades en la planificación de proyectos institucionales y una limitada capacidad para generar reportes confiables en el corto plazo.

Esta situación refleja la ausencia de un mecanismo tecnológico unificado que permita centralizar y automatizar la gestión de evidencias y de tiempos de jornada. Dicha carencia afecta directamente la eficiencia administrativa, la transparencia en la toma de decisiones y la capacidad institucional para responder a los procesos de acreditación y planificación académica.

En consecuencia, se identifica como problema la necesidad de implementar un sistema web integrado que atienda estas deficiencias, facilite el cumplimiento de estándares de calidad y fortalezca la gestión estratégica de la Universidad Nacional.
